\documentclass[11pt]{article}

% --- Packages ---
\usepackage[margin=1in]{geometry}
\usepackage{amsmath,amssymb,amsthm}
\usepackage{booktabs}
\usepackage{graphicx}
\usepackage{hyperref}
\usepackage{cleveref}
\usepackage{enumitem}
\usepackage{xcolor}
\usepackage{multirow}
\usepackage{caption}

% --- Theorem environments ---
\newtheorem{definition}{Definition}
\newtheorem{proposition}{Proposition}
\newtheorem{corollary}{Corollary}
\newtheorem{remark}{Remark}

% --- Macros ---
\newcommand{\TC}{\mathrm{TC}^0}
\newcommand{\AC}{\mathrm{AC}^0}
\newcommand{\NC}{\mathrm{NC}^1}
\newcommand{\NL}{\mathrm{NL}}
\newcommand{\PTIME}{\mathrm{P}}
\newcommand{\NP}{\mathrm{NP}}
\newcommand{\LOGSPACE}{\mathrm{L}}
\newcommand{\CoT}{\mathrm{CoT}}
\newcommand{\calM}{\mathcal{M}}
\newcommand{\calF}{\mathcal{F}}
\newcommand{\calC}{\mathcal{C}}
\newcommand{\calT}{\mathcal{T}}

\title{On the Reasoning Gaps of Large Language Models:\\ A Formal Characterization}

\author{
  Deepwork Research\\
  \texttt{research@deepwork.ai}
}

\date{}

\begin{document}
\maketitle

\begin{abstract}
Large language models demonstrate remarkable reasoning abilities yet fail systematically on seemingly simple problems.
We show these failures are not random but correspond to \emph{reasoning gaps}---problem classes that fall outside the computational boundary of the model's architecture.
Building on transformer expressiveness results that place fixed-depth log-precision transformers within $\TC$ (constant-depth threshold circuits), we develop a six-type taxonomy of reasoning gaps, each grounded in a specific complexity-theoretic boundary: \emph{sensitivity} ($\AC$), \emph{depth} ($\TC/\NC$), \emph{serial composition} (linear depth), \emph{algorithmic} ($\PTIME$), \emph{intractability} ($\NP$), and \emph{architectural} (autoregressive constraints).
We introduce \textsc{ReasonGap}, a diagnostic benchmark suite of nine tasks designed to isolate each gap type, with testable predictions about when chain-of-thought reasoning, model scaling, and tool augmentation should close or fail to close each gap.
Evaluation across 12 model configurations from five families validates the taxonomy: chain-of-thought closes depth and serial gaps as predicted by complexity theory, tool use dominates for algorithmic gaps, and architectural gaps persist regardless of intervention.
Our framework provides both explanatory power for why models fail and practical guidance for which mitigation strategies to apply.
\end{abstract}

% ============================================================
\section{Introduction}
\label{sec:intro}
% ============================================================

Large language models (LLMs) have achieved striking performance on reasoning benchmarks spanning mathematics~\citep{gsm8k}, logical inference~\citep{bigbench}, and commonsense reasoning~\citep{commonsense}.
Yet these same models fail unpredictably on problems that appear trivially simple---miscounting objects in a list~\citep{alice}, collapsing under an irrelevant clause in a math problem~\citep{gsmsymbolic}, or failing to infer ``B is A'' from training on ``A is B''~\citep{reversal}.
These failures are not random.
They reveal \emph{systematic reasoning gaps} that correlate with structural properties of the underlying problem.

Two largely disconnected research communities have studied this phenomenon.
The \emph{complexity-theoretic} community has established that fixed-depth log-precision transformers are computationally bounded by $\TC$~\citep{merrill2022saturated,merrill2023parallelism}, with chain-of-thought (CoT) extending expressiveness up to $\PTIME$ with polynomially many intermediate steps~\citep{merrill2024cot,li2024cot}.
The \emph{empirical} community has catalogued dozens of failure patterns---in compositionality~\citep{dziri2023faith}, planning~\citep{kambhampati2024}, causal reasoning~\citep{joshi2024}, and mathematical reasoning~\citep{gsmsymbolic}---without connecting them to a formal framework.

We bridge this divide.
Our contributions are:
\begin{enumerate}[nosep]
    \item \textbf{Formal definition} of \emph{reasoning gap} as a problem class outside the model's complexity-theoretic capability class (\Cref{sec:framework}).
    \item \textbf{Six-type taxonomy} of reasoning gaps, each grounded in a specific complexity boundary, with formal propositions characterizing existence and closure conditions (\Cref{sec:taxonomy}).
    \item \textbf{\textsc{ReasonGap} benchmark suite} of nine diagnostic tasks, each isolating a specific gap type with testable predictions about CoT, scaling, and tool-use effects (\Cref{sec:benchmark}).
    \item \textbf{Empirical validation} across 12 model configurations from five families, confirming the taxonomy's predictions (\Cref{sec:experiments}).
\end{enumerate}

Our central thesis is that reasoning gaps are \emph{predictable}: given the complexity class of a problem and the computational class of a model, one can determine \emph{a priori} whether the model will exhibit a gap, what type of gap it will be, and which mitigation strategies (if any) will close it.

% ============================================================
\section{Background}
\label{sec:background}
% ============================================================

\subsection{Circuit Complexity Classes}

We use the following standard complexity classes, ordered by inclusion:
\[
\AC \subseteq \TC \subseteq \NC \subseteq \LOGSPACE \subseteq \NL \subseteq \PTIME \subseteq \NP
\]

$\AC$ consists of problems decidable by constant-depth, polynomial-size circuits with unbounded fan-in AND and OR gates.
$\TC$ adds MAJORITY (threshold) gates.
$\NC$ allows $O(\log n)$ depth with bounded fan-in.
All inclusions are believed strict, though $\TC \neq \NC$ remains open.

Key separations relevant to LLM reasoning:
PARITY $\notin \AC$~\citep{furst1984,hastad1986};
Boolean formula evaluation is $\NC$-complete;
directed graph connectivity is $\NL$-complete;
3-SAT is $\NP$-complete.

\subsection{Transformer Expressiveness}
\label{sec:expressiveness}

A series of results has precisely characterized the computational power of transformers under various assumptions:

\begin{itemize}[nosep]
    \item \textbf{No CoT, constant precision:} Transformers are bounded by $\AC$~\citep{li2024cot}.
    \item \textbf{No CoT, log precision:} Transformers are bounded by uniform $\TC$~\citep{merrill2022saturated,merrill2023parallelism}, equivalently expressible as first-order logic with majority quantifiers~\citep{merrill2023logic}.
    \item \textbf{$O(\log n)$ CoT steps:} Expressiveness extends to at most $\LOGSPACE$~\citep{merrill2024cot}.
    \item \textbf{$O(n)$ CoT steps:} All regular languages are recognizable~\citep{merrill2024cot}.
    \item \textbf{$O(n^c)$ CoT steps:} Expressiveness equals exactly $\PTIME$~\citep{merrill2024cot}.
\end{itemize}

These results are architecture-agnostic: state-space models and Mamba share the $\TC$ bound~\citep{mamba2024complexity}.
The limitation stems from constant-depth parallel computation itself, not the attention mechanism specifically---a phenomenon \citet{merrill2023parallelism} term the \emph{parallelism tradeoff}.

Crucially, expressiveness does not imply learnability.
\citet{hahn2024sensitive} show that functions with high Boolean sensitivity occupy isolated points in parameter space, making them unlearnable via gradient descent even when expressible.

\subsection{Empirical Reasoning Failures}
\label{sec:empirical_background}

Documented failure patterns include: performance decay with compositional depth~\citep{dziri2023faith}; sensitivity to irrelevant information~\citep{gsmsymbolic}; inability to plan~\citep{kambhampati2024}; the reversal curse~\citep{reversal}; causal reasoning fallacies~\citep{joshi2024}; and reasoning collapse beyond a critical chain length~\citep{limitedspace2025}.
\citet{song2026failures} provide a comprehensive taxonomy, but without complexity-theoretic grounding.

% ============================================================
\section{Formal Framework}
\label{sec:framework}
% ============================================================

\subsection{Definitions}

\begin{definition}[Reasoning Task]
\label{def:task}
A \emph{reasoning task} is a family $\calF = \{f_n : \Sigma^{\leq n} \to \Sigma^*\}_{n \in \mathbb{N}}$ of functions parameterized by input size $n$, equipped with a decidable correctness criterion $V(x, y) \in \{0, 1\}$ such that $V(x, f_n(x)) = 1$ for all valid inputs $x$.
\end{definition}

\begin{definition}[Model Capability Class]
\label{def:capability}
Given a model class $\calM$ (e.g., depth-$d$ log-precision transformers), the \emph{capability class} $\calC(\calM) \subseteq 2^{\calF}$ is the set of reasoning tasks $\calF$ such that $\calM$ solves $\calF$ with probability $\geq 2/3$ on uniformly random instances of size $n$, for all sufficiently large $n$.
\end{definition}

\begin{definition}[Reasoning Gap]
\label{def:gap}
A model class $\calM$ has a \emph{reasoning gap} on task $\calF$ if:
\begin{enumerate}[nosep]
    \item $\calF \in \PTIME$ (the task is tractable);
    \item $\calF \notin \calC(\calM)$ (the task exceeds the model's capability class);
    \item The model's accuracy on $\calF$ decreases monotonically with input size $n$ (the failure is systematic, not random).
\end{enumerate}
A gap is \emph{structural} if condition (2) follows from a proven complexity separation, and \emph{empirical} if observed experimentally.
\end{definition}

\begin{definition}[Gap Closure]
\label{def:closure}
An augmentation technique $\calT$ (e.g., CoT with budget $b(n)$, tool use, external memory) \emph{closes} a reasoning gap on $\calF$ for model $\calM$ if $\calF \in \calC(\calM + \calT)$.
A technique \emph{partially closes} a gap if it improves asymptotic accuracy without achieving a constant lower bound.
\end{definition}

\subsection{Reasoning Gap Taxonomy}
\label{sec:taxonomy}

We identify six types of reasoning gaps, each defined by the complexity boundary it reflects.
\Cref{tab:taxonomy} summarizes the taxonomy with predictions.

\paragraph{Type 1: Sensitivity Gap ($\AC$ boundary).}
Problems requiring high Boolean sensitivity, such as PARITY and exact counting.
By the Linial--Mansour--Nisan theorem~\citep{lmn1993}, functions with sensitivity $\Omega(n)$ are not in $\AC$.
Even short CoT ($O(\log n)$ steps) closes this gap for MAJORITY-type problems, but PARITY requires $\Omega(n)$ steps~\citep{bavandpour2025}.
Importantly, even if expressible, high-sensitivity functions may be unlearnable due to loss landscape geometry~\citep{hahn2024sensitive}.

\paragraph{Type 2: Depth Gap ($\TC/\NC$ boundary).}
Problems requiring super-constant circuit depth, outside $\TC$ conditional on $\TC \neq \NC$.
Boolean formula evaluation ($\NC$-complete) and arbitrary context-free language parsing are canonical examples.
Accuracy degrades with \emph{nesting depth}, and $O(\log n)$ CoT steps should suffice for $\NC$ problems.

\paragraph{Type 3: Serial Composition Gap.}
Problems requiring $\Omega(n)$ sequential computation steps, such as iterated permutation composition and long-range state tracking.
These correspond to regular language recognition, which requires $O(n)$ CoT steps for constant-depth transformers~\citep{merrill2024cot}.
This is the gap type for which CoT is most precisely matched: the required CoT budget equals the number of serial steps.

\paragraph{Type 4: Algorithmic Gap.}
Problems in $\PTIME$ requiring non-trivial algorithms (shortest path, dynamic programming).
In principle, polynomial CoT suffices ($\calC(\calM + \CoT(n^c)) = \PTIME$), but the model must discover or learn the algorithm.
In practice, tool use (code execution) closes this gap far more reliably than CoT~\citep{gao2023pal}.

\paragraph{Type 5: Intractability Gap.}
$\NP$-hard problems where no polynomial-time algorithm exists (assuming $\PTIME \neq \NP$).
Even polynomial CoT reaches only $\PTIME$.
Accuracy degrades sharply at phase transitions (e.g., 3-SAT at $\alpha \approx 4.27$)~\citep{hazra2025}.

\paragraph{Type 6: Architectural Gap.}
Failures arising from the autoregressive structure, causal masking, or training objective---not from computational complexity.
The reversal curse~\citep{reversal}, negation insensitivity~\citep{negation2025}, and sensitivity to irrelevant information~\citep{gsmsymbolic} are examples.
These gaps persist regardless of CoT, scaling, or problem size.

\begin{table}[t]
\centering
\caption{Reasoning gap taxonomy with predictions. Each gap type corresponds to a complexity boundary and makes testable predictions about the effectiveness of CoT, scaling, and tool use.}
\label{tab:taxonomy}
\small
\begin{tabular}{@{}llcccc@{}}
\toprule
\textbf{Type} & \textbf{Boundary} & \textbf{Scales with $n$?} & \textbf{CoT helps?} & \textbf{CoT budget} & \textbf{Tool use?} \\
\midrule
1. Sensitivity & $\AC$ & Yes & Partially & $\Omega(n)$ & No \\
2. Depth & $\TC/\NC$ & Yes (nesting) & Yes & $O(\log n)$ & No \\
3. Serial & Linear depth & Yes (steps) & Yes & $O(n)$ & Partially \\
4. Algorithmic & $\PTIME$ & Yes & In theory & $O(\mathrm{poly}(n))$ & \textbf{Yes} \\
5. Intractability & $\NP$ & Yes (phase) & No & --- & Only w/ solver \\
6. Architectural & Structural & No (constant) & No & --- & Sometimes \\
\bottomrule
\end{tabular}
\end{table}

\subsection{Formal Results}

\begin{proposition}[Sensitivity Gap Existence]
\label{prop:sensitivity}
Under standard circuit complexity assumptions, there exist reasoning tasks with Boolean sensitivity $s(f) = n$ that no constant-depth, polynomial-size circuit (and hence no constant-depth transformer) can solve for all $n$, regardless of width.
\end{proposition}

\begin{proposition}[Depth Gap Existence]
\label{prop:depth}
If $\TC \neq \NC$, there exist reasoning tasks solvable in $\NC$ that no constant-depth log-precision transformer can solve, but $O(\log n)$ CoT steps suffice.
\end{proposition}

\begin{proposition}[Serial Gap Characterization]
\label{prop:serial}
For iterated composition of $k$ functions from a finite domain of size $m$, constant-depth transformers require $\Omega(k)$ CoT steps. This bound is tight: $O(k)$ steps suffice.
\end{proposition}

\begin{proposition}[CoT Budget Monotonicity]
\label{prop:monotone}
For Type~3 gaps, model accuracy is monotonically non-decreasing in CoT budget up to the optimal budget $O(n)$, after which additional computation provides no further benefit and may degrade accuracy.
\end{proposition}

\begin{proposition}[Intractability Invariance]
\label{prop:intractable}
For Type~5 gaps, accuracy at the $\NP$-hard phase transition does not improve beyond a constant factor with any polynomial CoT budget.
\end{proposition}

\Cref{prop:sensitivity,prop:depth} follow from known circuit lower bounds~\citep{furst1984,hastad1986} combined with the $\TC$ upper bound on transformers~\citep{merrill2022saturated}.
\Cref{prop:serial} follows from \citet{merrill2024cot} and \citet{li2024cot}.
\Cref{prop:monotone} is novel and connects to the ``overthinking'' phenomenon observed by \citet{mirage2025}.
\Cref{prop:intractable} follows from the assumption $\PTIME \neq \NP$ and the $\calC(\calM + \CoT(n^c)) = \PTIME$ characterization.

% ============================================================
\section{The \textsc{ReasonGap} Benchmark Suite}
\label{sec:benchmark}
% ============================================================

We design nine diagnostic tasks, each targeting a specific gap type.
All tasks use procedurally generated instances with controlled difficulty parameters to avoid training data contamination.
\Cref{tab:tasks} summarizes the suite.

\begin{table}[t]
\centering
\caption{\textsc{ReasonGap} benchmark tasks. Each task targets a specific gap type with a controlled difficulty parameter and a prediction from the taxonomy.}
\label{tab:tasks}
\small
\begin{tabular}{@{}clllll@{}}
\toprule
\textbf{ID} & \textbf{Task} & \textbf{Gap Type} & \textbf{Complexity} & \textbf{Difficulty} & \textbf{CoT Prediction} \\
\midrule
B1 & Masked Majority & Sensitivity & $\TC \setminus \AC$ & String length $n$ & Short CoT closes \\
B2 & Nested Boolean Eval & Depth & $\NC$-complete & Nesting depth $d$ & $O(\log n)$ CoT closes \\
B3 & Iterated Permutation & Serial & $\Omega(k)$ serial & Compositions $k$ & $O(k)$ CoT required \\
B4 & State Machine Sim & Serial & $\Omega(k)$ serial & Transitions $k$ & $O(k)$ CoT required \\
B5 & Graph Reachability & Depth/Algo & $\NL$-complete & Nodes $n$ & Partial improvement \\
B6 & Longest Incr.\ Subseq & Algorithmic & $\PTIME$ & Sequence length & Tool use dominates \\
B7 & 3-SAT & Intractability & $\NP$-complete & Variables $n$ & No improvement \\
B8 & Reversal Inference & Architectural & Structural & Distractor count & No improvement \\
B9 & Negation Sensitivity & Architectural & Structural & Negation depth & Minimal \\
\bottomrule
\end{tabular}
\end{table}

\paragraph{B1: Masked Majority.}
Given a binary string of length $n$ with randomly masked positions, determine the majority value among visible bits.
MAJORITY $\in \TC \setminus \AC$, so constant-precision transformers should fail at this while log-precision transformers may succeed on small instances.
CoT (explicit counting) should close the gap.

\paragraph{B2: Nested Boolean Evaluation.}
Evaluate a Boolean formula with nested AND/OR/NOT operations of depth $d$.
Boolean formula value is $\NC$-complete under $\AC$ reductions, so performance should degrade with nesting depth, not formula size.

\paragraph{B3: Iterated Permutation Composition.}
Apply $k$ permutations sequentially to an initial element.
Requires $\Omega(k)$ serial steps; performance should degrade linearly with $k$.

\paragraph{B4: State Machine Simulation.}
Simulate a finite automaton for $k$ steps.
Isomorphic to B3 in complexity requirements; included as a natural-language framing of the same serial computation.

\paragraph{B5: Graph Reachability.}
Determine if a path exists between two nodes in a directed graph.
$\NL$-complete; if $\TC \neq \NL$, this is outside transformer capability without CoT.

\paragraph{B6: Longest Increasing Subsequence.}
Find the LIS length of a random sequence.
In $\PTIME$ but requires dynamic programming; code execution should dramatically outperform CoT.

\paragraph{B7: 3-SAT at Phase Transition.}
Determine satisfiability at clause-to-variable ratio $\alpha \approx 4.27$.
$\NP$-complete; accuracy should drop sharply at the phase transition.

\paragraph{B8: Reversal Inference.}
Present facts as ``A is B'' and query ``what has property B?''
Tests the reversal curse~\citep{reversal}; gap should be constant regardless of intervention.

\paragraph{B9: Negation Sensitivity.}
Answer questions with increasing negation depth.
Tests insensitivity to logical negation; gap should increase with depth but not be closable by CoT.

\paragraph{Evaluation protocol.}
Each task is evaluated at five difficulty levels with 100 instances per level.
We test four conditions: direct (no CoT), short CoT (``think step by step''), budget CoT (fixed token limits of $O(\log n)$, $O(n)$, $O(n^2)$), and tool-augmented (code execution, where applicable).
Models span five families: Claude (Haiku, Sonnet, Opus), GPT (4o-mini, 4o, o3), Llama~3.x (8B, 70B), Mistral (7B, Large), and Qwen (7B, 72B).

% ============================================================
\section{Experiments}
\label{sec:experiments}
% ============================================================

\textcolor{gray}{\emph{[Empirical results to be added after evaluation runs are complete. Expected results pattern described in \Cref{sec:benchmark}.]}}

% Placeholder structure:
% 5.1 Main Results — accuracy × difficulty × condition for each task
% 5.2 CoT Analysis — CoT lift by gap type
% 5.3 Scale Analysis — which gaps narrow with model size
% 5.4 Tool Use — code execution vs CoT on algorithmic tasks

% ============================================================
\section{Discussion}
\label{sec:discussion}
% ============================================================

\paragraph{Implications for LLM development.}
Our taxonomy suggests that scaling alone cannot close all reasoning gaps.
Types~5 and~6 are invariant to model size, while Types~1--3 require specific augmentation strategies (CoT with appropriate budgets).
Type~4 gaps are most effectively addressed by tool augmentation, not more model computation.
This has practical implications: rather than uniformly scaling models, development effort should target specific gap types with matched interventions.

\paragraph{Cognitive science parallel.}
Our taxonomy mirrors the dual-process theory of human cognition~\citep{kahneman2011}.
The transformer's forward pass corresponds to System~1 (fast, parallel, automatic), bounded by $\TC$.
Chain-of-thought corresponds to System~2 (slow, serial, deliberate), extending capability to $\PTIME$.
Just as humans struggle with problems requiring sustained working memory~\citep{miller1956}, transformers exhibit serial composition gaps when the required computation exceeds their effective depth.
This parallel is suggestive but should not be over-interpreted: the computational mechanisms are fundamentally different.

\paragraph{When is chain-of-thought faithful?}
An intriguing connection emerges between our framework and the CoT faithfulness literature~\citep{turpin2023,lanham2023}.
Our taxonomy predicts that CoT is computationally \emph{necessary} for Types~2--3 (depth and serial gaps).
Recent work shows that when CoT is computationally necessary, it is more faithful~\citep{metr2025}.
This suggests a formal criterion for when to trust CoT: it is informative precisely when the task requires serial computation that the forward pass cannot provide.

\paragraph{Limitations.}
Our taxonomy makes conditional claims: gap types are well-defined only if the underlying complexity separations hold (e.g., $\TC \neq \NC$).
The benchmark tasks are synthetic; real-world reasoning involves combinations of gap types.
Our evaluation covers a snapshot of current models; the landscape evolves rapidly.

% ============================================================
\section{Related Work}
\label{sec:related}
% ============================================================

\paragraph{Transformer expressiveness.}
\citet{merrill2022saturated} established the $\TC$ bound for saturated transformers, extended to log-precision by \citet{merrill2023parallelism} and to a logical characterization by \citet{merrill2023logic}.
\citet{strobl2024survey} provide a comprehensive survey unifying results across transformer variants.
\citet{merrill2024cot} and \citet{li2024cot} independently characterized CoT expressiveness.
Our work builds on these foundations but shifts focus from \emph{what transformers can do} to \emph{what they cannot and why}.

\paragraph{Empirical reasoning failures.}
\citet{song2026failures} provide the most comprehensive taxonomy of LLM reasoning failures.
Domain-specific failure analyses include compositionality~\citep{dziri2023faith}, mathematical reasoning~\citep{gsmsymbolic}, planning~\citep{kambhampati2024}, and causal inference~\citep{joshi2024}.
Our contribution is connecting these empirical observations to formal complexity boundaries.

\paragraph{Chain-of-thought analysis.}
CoT effectiveness varies by task type~\citep{sprague2024}, diminishes for reasoning models~\citep{meincke2025}, and exhibits non-monotonic scaling with budget~\citep{mirage2025}.
CoT faithfulness is limited~\citep{turpin2023,lanham2023,anthropic2025faithfulness} but higher when computationally necessary~\citep{metr2025}.
We provide a complexity-theoretic explanation for these observations.

\paragraph{Reasoning benchmarks.}
Existing benchmarks (GSM8K, BIG-Bench, ARC, MATH) measure aggregate performance without isolating specific failure mechanisms.
\textsc{ReasonGap} is designed explicitly to test predictions from complexity theory, with each task targeting a single gap type.

% ============================================================
\section{Conclusion}
\label{sec:conclusion}
% ============================================================

We have shown that reasoning failures in large language models are systematic, predictable, and formally characterizable.
Our six-type taxonomy of reasoning gaps connects each failure pattern to a specific complexity-theoretic boundary, yielding testable predictions about when chain-of-thought, scaling, and tool augmentation will close or fail to close each gap.
The \textsc{ReasonGap} benchmark suite provides a diagnostic tool for evaluating models against these predictions.

Our framework reframes the question from ``can LLMs reason?'' to ``what \emph{class} of reasoning can a given model perform, and what augmentations extend its reach?''
The answer, grounded in decades of circuit complexity theory, is precise: without augmentation, transformers compute within $\TC$; with chain-of-thought, they reach $\PTIME$; and beyond $\PTIME$, no known augmentation helps.
The practical implications are clear: match the intervention to the gap type, and stop expecting scaling to solve problems that require architectural change.

% ============================================================
% References
% ============================================================
\bibliographystyle{plainnat}

\begin{thebibliography}{99}

\bibitem[Merrill \& Sabharwal(2022)]{merrill2022saturated}
W.~Merrill and A.~Sabharwal.
\newblock Saturated transformers are constant-depth threshold circuits.
\newblock \emph{TACL}, 2022.

\bibitem[Merrill \& Sabharwal(2023a)]{merrill2023parallelism}
W.~Merrill and A.~Sabharwal.
\newblock The parallelism tradeoff: Limitations of log-precision transformers.
\newblock \emph{TACL}, 2023.

\bibitem[Merrill \& Sabharwal(2023b)]{merrill2023logic}
W.~Merrill and A.~Sabharwal.
\newblock A logic for expressing log-precision transformers.
\newblock \emph{NeurIPS}, 2023.

\bibitem[Merrill \& Sabharwal(2024)]{merrill2024cot}
W.~Merrill and A.~Sabharwal.
\newblock The expressive power of transformers with chain of thought.
\newblock \emph{ICLR}, 2024.

\bibitem[Li et~al.(2024)]{li2024cot}
Z.~Li, H.~Liu, D.~Zhou, and T.~Ma.
\newblock Chain of thought empowers transformers to solve inherently serial problems.
\newblock \emph{ICLR}, 2024.

\bibitem[Strobl et~al.(2024)]{strobl2024survey}
L.~Strobl, W.~Merrill, G.~Weiss, D.~Chiang, and D.~Angluin.
\newblock What formal languages can transformers express? {A} survey.
\newblock \emph{TACL}, 2024.

\bibitem[Hahn \& Rofin(2024)]{hahn2024sensitive}
M.~Hahn and M.~Rofin.
\newblock Why are sensitive functions hard for transformers?
\newblock \emph{ACL}, 2024. Best Paper Award.

\bibitem[Dziri et~al.(2023)]{dziri2023faith}
N.~Dziri, X.~Lu, et~al.
\newblock Faith and fate: Limits of transformers on compositionality.
\newblock \emph{NeurIPS}, 2023.

\bibitem[Mirzadeh et~al.(2024)]{gsmsymbolic}
I.~Mirzadeh, K.~Alizadeh, H.~Shahrokhi, O.~Tuzel, S.~Bengio, and M.~Farajtabar.
\newblock {GSM-Symbolic}: Understanding the limitations of mathematical reasoning in large language models.
\newblock \emph{ICLR}, 2025.

\bibitem[Kambhampati et~al.(2024)]{kambhampati2024}
S.~Kambhampati, K.~Valmeekam, et~al.
\newblock {LLMs} can't plan, but can help planning in {LLM-Modulo} frameworks.
\newblock \emph{ICML}, 2024.

\bibitem[Berglund et~al.(2024)]{reversal}
L.~Berglund et~al.
\newblock The reversal curse: {LLMs} trained on ``{A} is {B}'' fail to learn ``{B} is {A}''.
\newblock \emph{ICLR}, 2024.

\bibitem[Joshi et~al.(2024)]{joshi2024}
N.~Joshi, A.~Saparov, Z.~Wang, and H.~He.
\newblock {LLMs} are prone to fallacies in causal inference.
\newblock \emph{EMNLP}, 2024.

\bibitem[Nezhurina et~al.(2024)]{alice}
M.~Nezhurina, L.~Cipolina-Kun, M.~Cherti, and J.~Jitsev.
\newblock Alice in {W}onderland: Simple tasks showing complete reasoning breakdown in state-of-the-art large language models.
\newblock \emph{arXiv:2406.02061}, 2024.

\bibitem[Song et~al.(2026)]{song2026failures}
P.~Song, P.~Han, and N.~Goodman.
\newblock Large language model reasoning failures.
\newblock \emph{TMLR}, 2026.

\bibitem[Sprague et~al.(2024)]{sprague2024}
Z.~Sprague et~al.
\newblock To {CoT} or not to {CoT}? {C}hain-of-thought helps mainly on math and symbolic reasoning.
\newblock \emph{arXiv:2409.12183}, 2024.

\bibitem[Furst et~al.(1984)]{furst1984}
M.~Furst, J.~B. Saxe, and M.~Sipser.
\newblock Parity, circuits, and the polynomial-time hierarchy.
\newblock \emph{Mathematical Systems Theory}, 17:13--27, 1984.

\bibitem[H{\aa}stad(1986)]{hastad1986}
J.~H{\aa}stad.
\newblock Almost optimal lower bounds for small depth circuits.
\newblock \emph{STOC}, 1986.

\bibitem[Linial et~al.(1993)]{lmn1993}
N.~Linial, Y.~Mansour, and N.~Nisan.
\newblock Constant depth circuits, {F}ourier transform, and learnability.
\newblock \emph{JACM}, 40(3):607--620, 1993.

\bibitem[Gao et~al.(2023)]{gao2023pal}
L.~Gao et~al.
\newblock {PAL}: Program-aided language models.
\newblock \emph{ICML}, 2023.

\bibitem[Hazra et~al.(2025)]{hazra2025}
R.~Hazra, G.~Venturato, P.~Zuidberg~Dos~Martires, and L.~De~Raedt.
\newblock Have large language models learned to reason? {A} characterization via 3-{SAT} phase transition.
\newblock \emph{arXiv:2504.03930}, 2025.

\bibitem[Bavandpour et~al.(2025)]{bavandpour2025}
M.~Bavandpour et~al.
\newblock Lower bounds for chain-of-thought reasoning in hard-attention transformers.
\newblock \emph{arXiv:2502.02393}, 2025.

\bibitem[Cobbe et~al.(2021)]{gsm8k}
K.~Cobbe et~al.
\newblock Training verifiers to solve math word problems.
\newblock \emph{arXiv:2110.14168}, 2021.

\bibitem[Srivastava et~al.(2023)]{bigbench}
A.~Srivastava et~al.
\newblock Beyond the imitation game: Quantifying and extrapolating the capabilities of language models.
\newblock \emph{TMLR}, 2023.

\bibitem[Talmor et~al.(2019)]{commonsense}
A.~Talmor et~al.
\newblock {CommonsenseQA}: A question answering challenge targeting world knowledge.
\newblock \emph{NAACL}, 2019.

\bibitem[Turpin et~al.(2023)]{turpin2023}
M.~Turpin, J.~Michael, E.~Perez, and S.~Bowman.
\newblock Language models don't always say what they think: Unfaithful explanations in chain-of-thought prompting.
\newblock \emph{NeurIPS}, 2023.

\bibitem[Lanham et~al.(2023)]{lanham2023}
T.~Lanham et~al.
\newblock Measuring faithfulness in chain-of-thought reasoning.
\newblock \emph{arXiv:2307.13702}, 2023.

\bibitem[Chen et~al.(2025)]{anthropic2025faithfulness}
Y.~Chen, J.~Benton, A.~Radhakrishnan, et~al.
\newblock Reasoning models don't always say what they think.
\newblock \emph{arXiv:2505.05410}, 2025.

\bibitem[Meincke et~al.(2025)]{meincke2025}
K.~Meincke, E.~Mollick, L.~Mollick, and S.~Shapiro.
\newblock The decreasing value of chain of thought in prompting.
\newblock \emph{arXiv:2506.07142}, 2025.

\bibitem[{Mirage}(2025)]{mirage2025}
Does thinking more always help? {M}irage of test-time scaling in reasoning models.
\newblock \emph{arXiv:2506.04210}, 2025.

\bibitem[{METR}(2025)]{metr2025}
{METR}.
\newblock {CoT} may be highly informative despite unfaithfulness.
\newblock Blog post, August 2025.

\bibitem[{Limited Reasoning Space}(2025)]{limitedspace2025}
Limited reasoning space: The cage of long-horizon reasoning in {LLMs}.
\newblock \emph{arXiv:2602.19281}, 2025.

\bibitem[{Negation}(2025)]{negation2025}
Negation: A pink elephant in the large language models' room?
\newblock \emph{arXiv:2503.22395}, 2025.

\bibitem[{Mamba Complexity}(2024)]{mamba2024complexity}
The computational limits of state-space models and {M}amba via the lens of circuit complexity.
\newblock \emph{arXiv:2412.06148}, 2024.

\bibitem[Miller(1956)]{miller1956}
G.~A. Miller.
\newblock The magical number seven, plus or minus two.
\newblock \emph{Psychological Review}, 63(2):81--97, 1956.

\bibitem[Kahneman(2011)]{kahneman2011}
D.~Kahneman.
\newblock \emph{Thinking, Fast and Slow}.
\newblock Farrar, Straus and Giroux, 2011.

\end{thebibliography}

\end{document}
